Show the convergence of the Euler product both absolute and non-absolute
and its generalized form

First note that the converge of a product \( \prod_{i=1}^\infty a_i \) is defined such that
the partial products \( \prod_{i=1}^N a_i \) are non-zero and converge. This holds iff.
\( \sum_{i=1}^\infty \log(a_i) \) converges.

By Euler’s product formula for s > 1
\(\zeta(s) = \prod_{p} \left(1 - \frac{1}{p^s}\right)^{-1}.\)
and its generalization
\(\zeta_K(s) = \prod_{\mathfrak{p}} \left(1 - \frac{1}{N(\mathfrak{p})^s}\right)^{-1}.\)
Note the power series expansion
\( \log(1 + z) = \sum^\infty_{n=1} = \frac{(-1)^n+1z^n}{n} = z - \frac{z^2}{2} + \frac{z^3}{3} + \dots \)
Then (the minus cancels with the minus inside the log in the expansion)
\( -\log(1 - \frac{1}{p^s}) = \sum_{n=1}^\infty \frac{1}{np^{ns}} \), which is conv. since 
\( |\frac{1}{p^s} = \frac{1}{p^\delta} < \frac{1}{p} < 1\)
and absolute convergence by
\( \sum_{p} |-\log(1 - \frac{1}{p^s})| = \sum_{p} |\frac{1}{np^{ns}}| \leq \sum_{p}\sum_{n=1}^\infty (|\frac{1}{p^{\delta}}|)^n\)
and from \( \frac{1}{p^\delta -1} < \frac{2}{p^\delta} \) (as \( \delta > 1 \))
\( \leq 2\sum_p \frac{1}{p^\delta} \leq 2 \sum_{n=1}^\infty \frac{1}{n^\delta}\), which converges as a Riemann series.

\[
\log \zeta_K(s) = \sum_{\mathfrak{p}} \sum_{m=1}^\infty \frac{1}{m N(\mathfrak{p})^{ms}}. \quad (9.18)
\]

For \( \zeta_K(s) \), take any \( \mathfrak{p} \) of \( \mathcal{O}_K \), let \( f_\mathfrak{p} \) denote its absolute inertia degree
then \( \mathcal{N}(\mathfrak{p}) = p^{f_\mathfrak{p}} \), where \( p \) is the rational prime lying below p. 

We rewrite again
\[
\log \zeta_K(s) = \sum_{\mathfrak{p}} \sum_{m=1}^\infty \frac{1}{m N(\mathfrak{p})^{ms}}.
\]
\[
= \sum_{p} \sum_{m=1}^\infty \sum_{\mathfrak{p}|p} \frac{1}{m N(\mathfrak{p})^{ms}}.
\]
now \( \frac{1}{m N(\mathfrak{p})^{ms}} = \frac{1}{mp^{f_\mathfrak{p}})^{ms}} \leq \frac{1}{mp^{ms}} \) 
and there are up to \( n = [K:\mathbb{Q}] \) prime ideals \( \mathfrak{p} \) which divide \( p \) by the class formula
so overall the term bounds to
\[
    n\sum_{p} \sum_{m=1}^\infty \frac{1}{mp^{ms}}
\]
which is absolutely conv. as before.



Define characters of groups and prove that their cardinality matches the order of the group

Let \( G \) be a finite abelian group. By a character \( \chi \) of \( G \), we mean a homomorphism \( \chi : G \rightarrow \{ z : |z| = 1, z \in \mathbb{C} \} \). If \( G \) has order \( n \) and \( e \) is the identity of \( G \), then for \( x \in G \), \( x^n = e \) which implies that \( \chi(x)^n = \chi(x^n) = \chi(e) = 1 \). So \( \chi(G) \) is contained in the group of \( n \)-th root of unity.

Observe that the set of all characters of a finite abelian group is a group under usual multiplication of mappings. 
This group \( \hat{G} = \text{Hom}(G, \mathbb{C}^{\times}) \)is called the character group of a finite abelian group \( G \).

Note that if \( G \) is a cyclic group of order \( m \) generated by \( a \), then \( G \) has exactly \( m \) characters \( \chi_0, \ldots, \chi_{m-1} \), defined by \( \chi_i(a) = \varepsilon^i \), where \( \varepsilon \) is a fixed primitive \( m \)-th root of unity. In fact this is true for every finite abelian group as asserted by the following proposition:

The number of characters of a finite abelian group equals the order of the group.

\( \textbf{Proof} \) We can write \( G = \prod_{i=1}^{s} G_i \) as a direct product of cyclic groups. 
Suppose \( G_i \) is generated by \( a_i \) and \( G_i \) has order \( m_i \). Any element \( x \in G \) is of the form
\[
x = a_1^{k_1} a_2^{k_2} \cdots a_s^{k_s}.
\]

So a character of \( G \) is completely determined if we define \( \chi(a_1), \ldots, \chi(a_s) \). Since \( a_i^{m_i} \) is the identity of \( G \), 
so \( \chi(a_i) \) is an \( (m_i) \)-th root of unity. 
Conversely let \( \varepsilon_i \) be an \( (m_i) \)-th root of unity, then for any \( x \in G \) given by above, we can define

\[
\chi(x) = \varepsilon_1^{k_1} \varepsilon_2^{k_2} \cdots \varepsilon_s^{k_s}
\] 

and this is well defined because the right hand side of (10.2) is independent of the choice of \( k_i \) (which are unique modulo \( m_i \)). 
Each root \( \varepsilon_i \) can be chosen in \( m_i \) ways. So there are \( m_1 m_2 \cdots m_s \) characters of \( G \). 


Show that characters of subgroups can be extended to the whole group and the existence of non-trivial characters:

Proposition. If \( G \) is a finite abelian group and \( H \) is a subgroup, then any character of \( H \) can be extended 
to a character of \( G \) and the number of such extensions equals the index of \( H \) in \( G \).

Proof. Let \( G \) be of order \( n \) and \( H \) be of order \( m \). When we restrict a character \( \chi \) of \( G \) to \( H \), 
we obtain a character of \( H \). 
It is clear that the map \( \Lambda \) sending \( \chi \) to the restriction map \( \chi|_H \) is a homomorphism 
from the character group \(\hat{G}\) of \( G \) into the character group \(\hat{H}\) of \( H \). 
Note that \( \chi \in \ker(\Lambda) \) if and only if \( \chi(g) = 1 \) for all \( g \in H \). 
It can be easily verified that the group \(\ker(\Lambda)\) is in one-to-one correspondence with the group of characters of \( G/H \). 
Thus by the previous proposition \(|\ker(\Lambda)| = |G/H| = \frac{n}{m}\). 
By the first isomorphism theorem of groups, \(\hat{G}/\ker(\Lambda) \cong \hat{H}\). 
So order of \(\Lambda(\hat{G}) = m\). But \(\hat{H}\) has order \( m \). Thus \(\Lambda(\hat{G}) = \hat{H}\) which implies that \(\Lambda\) is onto. 
Therefore, given any character of \( H \), it can be extended to a character of \( G \) 
and the number of such extensions equals \(\frac{n}{m}\).

Corollary. Let \( G \) be a finite abelian group. 
If \( g \neq e, g \in G \), then there exists a character \( \chi \) of \( G \) such that \( \chi(g) \neq 1 \).

Proof. Let \( H \) be the subgroup of \( G \) generated by \( g \), then \( H \neq \{e\} \). 
Let \(\psi\) be a non-trivial character of \( H \), so that \(\psi(g) \neq 1\). 
By the previous Proposition, we can extend the character \(\psi\) of \( H \) to a character \(\chi\) of \( G \). 
So \(\chi(g) = \psi(g) \neq 1\).

State and prove the identities for sums of characters

Proposition.
Let \( G \) be a finite abelian group of order \( n \) with identity \( e \). 
Let \(\hat{G}\) be the character group of \( G \) with identity \(\chi_0\). Then the following hold.

(i) For \( \chi \in \hat{G}, \quad \sum_{g \in G} \chi(g) \) equals \( n \) if \( \chi = \chi_0 \) and 0 otherwise.
(ii) For \( g \in G, \quad \sum_{\chi \in \hat{G}} \chi(g) \) equals \( n \) if \( g = e \) and 0 otherwise.

Proof. 

(i) If \( \chi = \chi_0 \), then \(\sum_{g \in G} \chi(g) = n\). Suppose \( \chi = \chi_0 \). Therefore there exists \( z \in G \) such that \( \chi(z) \neq 1 \). As \(\chi\) runs over all the elements of \( G \) so does \( g z \). Thus 
\[
\sum_{g \in G} \chi(g) = \sum_{g \in G} \chi(g z) = \chi(z) \sum_{g \in G} \chi(g).
\]

implying \( \sum_{g \in G} \chi(g) = 0 \) by assumption on \( \chi(z) \).
(ii) If \( g = e \), then clearly \(\sum_{\chi \in \hat{G}} \chi(x) = |\hat{G}| = n\). If \( g \neq e \), then by a Corollary, there exists a character 
\(\chi'\) of \( G \) such that \(\chi'(g) \neq 1\). As \( \chi \) runs over all the elements of \( \hat{G} \), so does \( \chi\chi' \).
Thus \(\sum_{\chi \in \hat{G}} \chi(g) = \sum_{\chi \in \hat{G}} \chi\chi'(g) = \chi'(g) \sum_{\chi \in \hat{G}} \chi(g).\) Since \(\chi'(g) \neq 1,\) 
we have \(\sum_{x \in \hat{G}} \chi(x) = 0.\) \(\quad \Box\)


Define Dirichlet characters

Notation. For any natural number \(m,\) let \(G_m\) denote the multiplicative group of all \\
reduced residue classes modulo \(m,\) i.e., those cosets of the group \(\mathbb{Z}/m\mathbb{Z}\) which are of \\
the form \(m\mathbb{Z} + a\) with \((a, m) = 1;\) \(G_m\) is a group of order \(\phi(m).\) The residue class \\
modulo \(n\) which contains an integer \(a\) will be denoted by \(\tilde{a}.\) To every character \(\chi\) of \\
\(G_m,\) one can associate a function \(\chi^*: \mathbb{Z} \longrightarrow \mathbb{C}\) defined by

\[
\chi^*(a) = 
\begin{cases}
    \chi(\tilde{a}) & \text{if } (a, m) = 1, \\
    0 & \text{if } (a, m) > 1.
\end{cases}
\]

Such a function \(\chi^*\) defined on \(\mathbb{Z}\) is called a numerical character modulo \(m\) or Dirichlet character. 
So a function \(\chi^*: \mathbb{Z} \longrightarrow \mathbb{C}\) is called a numerical character modulo \(m\) if it has the
following properties:

(i) \(\chi^*(ab) = \chi^*(a)\chi^*(b)\) for all \(a, b \in \mathbb{Z},\)
(ii) \(\chi^*(a) = 0\) if and only if \((a, m) > 1,\)
(iii) \(\chi^*(a) = \chi^*(a')\) if \(a \equiv a' \pmod{m}.\)

In view of a Proposition, the number of numerical characters modulo \(m\) is \(\phi(m).\) \\
In future, we shall denote the numerical character \(\chi^*\) corresponding to a character \(\chi\) \\
of \(G_m\) by \(\chi\) again. Corresponding to the trivial character \(\chi_0\) of \(G_m,\) we shall denote \\
by \(\chi_0,\) the numerical character modulo \(m\) which is defined by

\[
\chi_0(a) = 
\begin{cases}
    1 & \text{if } (a, m) = 1, \\
    0 & \text{if } (a, m) > 1.
\end{cases}
\]

It will be called the principal character modulo \(m.\).


State and prove the Pontryagin duality for the case of characters
and use it to determine sums of characters. 

Let \( \hat{G} \) be the abelian group of characters, it holds without proof 
for any subgroup \( H \subset G \) with SES \( 1 \longrightarrow H \longrightarrow G \longrightarrow G/H \longrightarrow 1 \)
also that \( 1 \longrightarrow \hat{G}/\hat{H} \longrightarrow \hat{G} \longrightarrow {H} \longrightarrow 1 \) is a SES.

The Pontryagin duality states \( G \simeq \hat{\hat{G}} \) via \(\psi(g)(\chi) = \chi(g) \).
Since \( |G| = |\hat{G}| = |\hat{\hat{G}}| \) it suffices to show that the map is injective, that is 
\( \psi(g)(\chi) = 1 \) for all \( \chi \in \hat{G} \) iff. \( g = 1 \).
For \( g \neq 1 \) we can construct  a non-trivial character \( \chi' \) on the subgroup \( \langle g \rangle \) generated by \( g \)
and extend it to all of G (extension property of characters), so \( \psi(g)(\chi') \neq 1 \).

Remark: There is a powerful generalized version stating
There is a canonical isomorphism \( G \simeq \hat{\hat{G}}\) between any locally compact abelian group G
and its double dual.

It follows the Lemma 13.8:
Let \( g \in G \), then \( \sum_{\chi \in \hat{G}} \chi(g) = |G| \) if \( g = 1 \) and otherwise \( 0 \).
Apply now the previous result for the \( \sum_{g \in G} \chi(g) \) for a fixed character now to \( \hat{G} \) instead of \( G \)
with the help of the Pontryagin duality.


Define \( L \)-series and prove their convergence.

For a numerical character \(\chi\) modulo \(m,\) we denote the sum of the series \(\sum_{n=1}^{\infty} \frac{\chi(n)}{n^s}\) by \(L(s, \chi)\) 
whenever it converges and its sum function is called \(L\)-function attached with character \(\chi.\) 
Applying a Proposition (with \(a_n = \frac{\chi(n)}{n^s}\)), we see that if we would have convergence for \(s > 1\)

\[
\sum_{n=1}^{\infty} \frac{\chi(n)}{n^s} = \prod_{p}\left(1 - \frac{\chi(p)}{p^s}\right)^{-1},
\]

where the product in the above equation is taken over all rational primes \(p.\)


Proposition
Let \(\chi \neq \chi_0\) be a numerical character modulo \(m,\) then the series \(\sum_{n=1}^{\infty} \frac{\chi(n)}{n^s}\) 
converges uniformly and absolute on all compact subsets of \(K \subset \mathbb{C}\) with \(k \in K \implies \mathcal{Re}(k) > 1 \) 
and if \(\chi \neq \chi_0\) then it converges also for \(\mathcal{Re}(s)\) in \((0, \infty).\) 

Proof. 
Since \( |\chi(n)| \leq 1 \) this series is bounded by \( \zeta(n) \), which converges absolutely under the given conditions.

Remark:
In particular the Euler product induced by \( L(s, \chi) \) converges absolutely as well using the fact that \( \chi \) is strictly multiplicative.



State and prove the simplification of Dirichlet's class number for cyclotomic fields

Theorem
Let \(K = \mathbb{Q}(\zeta)\) where \(\zeta\) is a primitive \(m\)th root of unity, \(m \geq 3\). Let \(R\) denote the regulator of \(K\), then the class number \(h\) of \(K\) is given by
\[
h = \frac{w \sqrt{|d_K|}}{(2\pi)^{\frac{\phi(m)}{2}} R} F(1) \prod_{\chi \neq \chi_0} L(1, \chi).
\]

Proof. 
Keeping in mind that \( K \) has no real isomorphism and Dirichlet’s Class Number formula, we have
\[
\lim_{s \to 1^+} (s - 1) \zeta_K(s) = h \frac{(2\pi)^{\frac{\phi(m)}{2}} R}{w \sqrt{|d_K|}},
\]

where \( h \) is the class number of \( K \), \( R \) is the regulator of \( K \) and \( w \) is the number of roots of unity contained in \( K \). 
By Euler’s Product formula, \(\zeta_K(s) = \prod_{\mathfrak{p}}\left(1 - \frac{1}{N(\mathfrak{p})^s}\right)^{-1}\) for \( s > 1 \),
where \(\mathfrak{p}\) runs over all non-zero prime ideals of \(\mathcal{O}_K\). We shall first write the above product in a clear manner.

For \(s > 0\), define \(G(s) = \prod_{p \mid m} \prod_{\mathfrak{p} \mid p} \left(1 - \frac{1}{N(\mathfrak{p})^s}\right)^{-1}\). 
If \(\mathfrak{p}\) is a prime ideal of \(\mathcal{O}_K\) lying over a rational prime \(p\) not dividing \(m\), 
then we denote the absolute residual degree of \(\mathfrak{p}\) by \(f_p\), 
it is independent of the choice of \(\mathfrak{p}\) lying over \(p\). In fact \(f_p\) is the smallest positive integer such that 
\(p^{f_p} \equiv 1 \pmod{m}\) by a previous Theorem. Also the number of distinct prime ideals \(\mathfrak{p}\) of \(\mathcal{O}_K\) lying over \(p\) is 
\(\frac{\phi(m)}{f_p}\). Thus for \(s > 1\), we can write
\[
\zeta_K(s) = G(s) \prod_{p \nmid m} \left(1 - \frac{1}{p^{sf_p}}\right)^{-\frac{\phi(m)}{f_p}}.
\]

Consider now the expression \( \prod_{\chi} L(\chi, s) \), where the product is taken over all Dirichlet characters modulo \( m \).
We can rewrite it to \(\prod_{\chi} (1 - \chi(p) p^{-s})^{-1}\). For \(p \mid m\) this is 1. So let \(p \nmid m\). 
By chap., (10.3), \(f\) is the order of \(p\) mod \(m\) in \((\mathbb{Z}/m\mathbb{Z})^\ast\) and \(e = 1\). 
Since \(efr = \varphi(m)\), the quotient \(r = \varphi(m)/f\) is the index of the subgroup \(G_p\) generated by \(p\) in 
\(G = (\mathbb{Z}/m\mathbb{Z})^\ast\). 
Associating \(\chi \mapsto \chi(p)\) defines an isomorphism \(\widehat{G}_p \cong \mu_f\), and gives the exact sequence
\[
    1 \longrightarrow \widehat{G}/\widehat{G}_p \longrightarrow \widehat{G} \longrightarrow \widehat{G}_p \longrightarrow 1,
\]
We therefore find \(r = \#(\widehat{G}/\widehat{G}_p) = (G : G_p)\) elements in the preimage of \(\chi(p)\). It follows that
\[
\prod_{\chi} (1 - \chi(p) p^{-s})^{-1} = \prod_{\zeta \in \mu_f} (1 - \zeta p^{-s})^{-r} = (1 - p^{-fs})^{-r}.
\]
Inserting this into our previous equations and keeping in mind the derivation, yields
\[
\zeta_K(s) = G(s) \prod_\chi \prod_p \left(1 - \frac{\chi(p)}{p^s}\right)^{-1} = G(s) \prod_\chi \prod_p L(s, \chi),
\]

If \(\chi = \chi_0\) is the principal character modulo \(m\), then for \(s > 1\)
\[
L(s, \chi_0) = \prod_p \left(1 - \frac{\chi_0(p)}{p^s}\right)^{-1} = \prod_{p \nmid m} \left(1 - \frac{1}{p^s}\right)^{-1} = \zeta(s) \prod_{p \mid m} \left(1 - \frac{1}{p^s}\right),
\]
because \(\zeta(s) = \prod_p \left(1 - \frac{1}{p^s}\right)^{-1}\) by Euler’s product formula. Thus we have shown that for \(s > 1\)
\[
\zeta_K(s) = F(s) \zeta(s) \prod_{\chi \neq \chi_0} L(s, \chi),
\]
where
\[
F(s) = \prod_{p \mid m} \left(1 - \frac{1}{p^s}\right) \prod_{p \mid m} \prod_{\mathfrak{p} \mid p} \left(1 - \frac{1}{N(\mathfrak{p})^s}\right)^{-1}.
\]

On multiplying both sides of (10.7) by \(s - 1\) and taking limit as \(s \to 1^+\), it follows from Dirichlet's class formula
and the convergence of the \( L \)-series,

\[
    h\kappa = F(1) \prod_{\chi \neq \chi_0} L(1, \chi), \quad \text{where} \quad \kappa = \frac{(2\pi)^{\frac{\phi(m)}{2}} R}{w \sqrt{|d_K|}},
\]
with \(R\) the regulator \(K\), \(w\) the number of roots of unity in \(K\) and \(L(1, \chi) = \sum_{n=1}^{\infty} \frac{\chi(n)}{n}\).
It has been shown that \( w = m \) or \( 2m \) according as \( m \) is even or odd, which finishes the proof.


Derive from the simplified class formula of cyclotomic fields, that \(L(1, \chi)\) is non-zero.

Corollary (directly from simplified class formula)
Let \(K = \mathbb{Q}(\zeta)\), where \(\zeta\) is a primitive \(p\)th root of unity, \(p\) an odd prime. 
Then the class number \(h_p\) of \(K\) is given by 
\[
h_p = \frac{2p^{(p/2)}}{(2\pi)^{p-1} R} \prod_{\chi \neq \chi_0} L(1, \chi),
\]
where \(R\) is the regulator of \(K\) and \(\chi\) runs over all non-principal numerical characters modulo \(p\).
Therefore if \(\chi\) is a non-principal numerical character modulo \(m\), then \(L(1, \chi)\) is non-zero.


State and prove Dirichlet's Prime Number Theorem

Theorem.
Every arithmetic progression \( \{a + km\}_{k \in \mathbb{N}}, (a, m) = 1  \)
contains infinitely many prime numbers.

Proof. 
Let \(\chi\) be a Dirichlet character \(\mod \, m\). Then one has, for \(\text{Re}(s) > 1\),
\[
\log L(\chi, s) = -\sum_p \log(1 - \chi(p)p^{-s}) = \sum_p \sum_{m=1}^{\infty} \frac{\chi(p^m)}{mp^{ms}} 
\]
\[
= \sum_p \frac{\chi(p)}{p^{s}} \sum_p \sum_{n=2}^{\infty} \frac{\chi(p^n)}{np^{ns}} = \sum_p \frac{\chi(p)}{p^s} + g_{\chi}(s),
\]
where \(g_{\chi}(s)\) is holomorphic for \(\text{Re}(s) > \frac{1}{2}\).
Indeed

\[
 \sum_p \sum_{n=2}^{\infty} \frac{\chi(p^n)}{np^{ns}} \sum_p \sum_{n \geq 2} \frac{1}{np^{ns}}
\]
\[
  \leq \sum_p \sum_{n=2}^{\infty} \frac{1}{p^{ns}} \leq \sum_p \sum_{n=2}^{\infty} \frac{1}{p^n} 
\]
\[
 = \sum_p \frac{1}{p(p-1)} \leq \zeta(2).
\]


Multiplying by \(\chi(a^{-1})\) and summing over all characters \(\mod \, m\) yields
\[
\sum_{\chi} \chi(a^{-1}) \log L(\chi, s) = \sum_{\chi} \sum_p \frac{\chi(a^{-1}p)}{p^s} + g(s)
\]
\[
= \sum_{b=1}^{m} \sum_{\chi} \chi(a^{-1}b) \sum_{p \equiv b (m)} \frac{1}{p^s} + g(s)
\]
\[
= \sum_{p \equiv a (m)} \frac{\varphi(m)}{p^s} + g(s).
\]

Note here that
\[
\sum_{\chi} \chi(a^{-1}b) = 
\begin{cases} 
0, & \text{if } a \neq b, \\
\varphi(m) = \#(\mathbb{Z}/m\mathbb{Z})^{*}, & \text{if } a = b.
\end{cases}
\]

When we pass to the limit \( s \to 1 \) (\(s\) real \(> 1\)), \(\log L(\chi, s)\) stays bounded for \(\chi \neq \chi^0\) because \(L(\chi, 1) \neq 0\), whereas \(\log L(\chi^0, s) = \sum_{p \nmid m} \log(1 - p^{-s}) + \log \zeta(s) \) tends to \( \infty \) because \(\zeta(s)\) has a pole. 
The left-hand side of the above equation therefore tends to \(\infty\) and since \(g(s)\) is holomorphic at \(s = 1\) we find
\[
\lim_{s \to 1} \sum_{p \equiv a (m)} \frac{\varphi(m)}{p^s} = \infty.
\]

Thus the sum cannot consist of only finitely many terms, and the theorem is proved.





\[ 
    \sum_{n=2}^\infty |\frac{\chi(p^n)}{p^{ns}}| \leq \sum_{n=2}^\infty \frac{1}{p^{n\delta}} = \frac{1}{p^{\delta}} \frac{1}{1 - \frac{1}{p^{\delta}}} \\
    = \frac{1}{p^{\delta}} \frac{1}{1 - \frac{1}{p^{\delta}}} < \frac{4}{p^{2\delta}}
\]
and \( p \geq 2, \sigma > \frac{1}{2} \).

